% Gerado por docs/triagem/sintese/scripts/gera_tex.py. Não editar à mão.
\chapter{Estudos primários incluídos nas revisões}
\label{ap:estudos}

Lista gerada por \textit{script} a partir dos metadados de cada estudo (autores pelo DOI, no OpenAlex). O identificador (ID) é o usado no texto. SB indica estudo incluído pelo \textit{snowballing}. A nota de qualidade vai de 0 a 11.

\section{RSL 1}
\begin{small}
\begin{longtable}{lp{11.2cm}r}
\hline
\textbf{ID} & \textbf{Referência} & \textbf{Nota} \\ \hline\endhead
RSL1-0183 & Müter et al. (2024). Enhancing Open-Source Intelligence: Introducing the Functional Intelligence Tool (FINT). \textit{2024 9th International Conference on Big Data Analytics, ICBDA 2024}. DOI: \url{10.1109/ICBDA61153.2024.10607337}. & 7 \\
RSL1-0211 & Konstantakos et al. (2023). Self-supervised visual learning for analyzing firearms trafficking activities on the Web. \textit{Proceedings - 2023 IEEE International Conference on Big Data, BigData 2023}. DOI: \url{10.1109/BigData59044.2023.10386795}. & 9 \\
RSL1-0214 & Ngo e Le?Khac (2023). Ontology-based case study management towards bridging training and actual investigation gaps in digital forensics. \textit{Forensic Science International: Digital Investigation}. DOI: \url{10.1016/j.fsidi.2023.301621}. & 6 \\
RSL1-0288 & Ndlovu et al. (2023). A Situational Awareness Tool using Open-Source Intelligence (OSINT) and Artificial Intelligence (AI). \textit{Proceedings - 2023 IEEE International Conference on Advances in Data-Driven Analytics and Intelligent Systems, ADACIS 2023}. DOI: \url{10.1109/ADACIS59737.2023.10424256}. & 9 \\
RSL1-0309 & Alif et al. (2025). Advancing Dark Web Data Extraction: An Extensible Framework for Multimodal Threat Intelligence. \textit{2025 28th International Conference on Computer and Information Technology, ICCIT 2025}. DOI: \url{10.1109/ICCIT68739.2025.11491194}. & 5 \\
RSL1-0340 & Abate et al. (2023). Artificial intelligence to fight illicit trafficking of cultural property. \textit{International Archives of the Photogrammetry, Remote Sensing and Spatial Information Sciences - ISPRS Archives}. DOI: \url{10.5194/isprs-Archives-XLVIII-M-2-2023-3-2023}. & 8 \\
RSL1-0344 & Wan et al. (2024). Link Analysis of Conflict through Signed Spectral Embedding: The Russo-Ukrainian War. \textit{2024 International Conference on Military Communication and Information Systems, ICMCIS 2024}. DOI: \url{10.1109/ICMCIS61231.2024.10540798}. & 9 \\
RSL1-0348 & Biagio et al. (2024). MARPLE: A Framework for Social Media Threat Intelligence. \textit{International Conference on Artificial Intelligence, Computer, Data Sciences, and Applications, ACDSA 2024}. DOI: \url{10.1109/ACDSA59508.2024.10467738}. & 10 \\
RSL1-0368 & Pala et al. (2025). An Enhanced Framework for Surveillance: Integrating OSINT, SOCMINT, AI-Driven Analysis, and Digital HUMINT against Terrorism and Exploitation Networks. \textit{INDISCON 2025 - IEEE 6th India Council International Subsections Conference, Proceedings}. DOI: \url{10.1109/INDISCON66021.2025.11251813}. & 5 \\
RSL1-0381 & Margaros et al. (2026). Detecting FIMI: A methodological framework for testing OSINT tools in TTPs detection. \textit{Open Research Europe}. DOI: \url{10.12688/openreseurope.21702.2}. & 9 \\
RSL1-0466 & Huang et al. (2025). A Blockchain-Based Framework for OSINT Evidence Collection and Identification. \textit{Future Internet}. DOI: \url{10.3390/fi17120551}. & 7 \\
RSL1-0494 & Nitzl et al. (2025). The use of artificial intelligence in military intelligence: an experimental investigation of added value in the analysis process. \textit{Frontiers in Human Dynamics}. DOI: \url{10.3389/fhumd.2025.1540450}. & 10 \\
RSL1-0561 & Boom e Veenman (2018). Finding Dutch natives in online forums. \textit{Forensic Sciences Research}. DOI: \url{10.1080/20961790.2018.1482042}. & 10 \\
RSL1-0612 & Chae et al. (2019). A system approach for evaluating current and emerging army open-source intelligence tools. \textit{SysCon 2019 - 13th Annual IEEE International Systems Conference, Proceedings}. DOI: \url{10.1109/SYSCON.2019.8836885}. & 6 \\
RSL1-0618 & Johnsen e Franke (2019). The impact of preprocessing in natural language for open source intelligence and criminal investigation. \textit{Proceedings - 2019 IEEE International Conference on Big Data, Big Data 2019}. DOI: \url{10.1109/BigData47090.2019.9006006}. & 10 \\
RSL1-0662 & Shen e Chow (2020). Time and location topic model for analyzing lihkg forum data. \textit{Proceedings - 2020 13th Systematic Approaches to Digital Forensic Engineering, SADFE 2020}. DOI: \url{10.1109/SADFE51007.2020.00009}. & 8 \\
RSL1-0779 & Maybir e Chapman (2021). Web scraping of ecstasy user reports as a novel tool for detecting drug market trends. \textit{Forensic Science International: Digital Investigation}. DOI: \url{10.1016/j.fsidi.2021.301172}. & 9 \\
RSL1-0816 & Wijnberg e Le?Khac (2021). Identifying interception possibilities for WhatsApp communication. \textit{Forensic Science International: Digital Investigation}. DOI: \url{10.1016/j.fsidi.2021.301132}. & 8 \\
RSL1-0854 & Williams e Stephens (2020). Analysis of the ‘Open Source Internet Research Tool’: A Usage Perspective from UK Law Enforcement. \textit{IFIP Advances in Information and Communication Technology}. DOI: \url{10.1007/978-3-030-57404-8_26}. & 8 \\
RSL1-0856 & Zhao et al. (2015). OSIA: Open source intelligence analysis system based on cloud computing and domestic platform. \textit{Proceedings - 2015 2nd International Conference on Information Science and Control Engineering, ICISCE 2015}. DOI: \url{10.1109/ICISCE.2015.89}. & 4 \\
RSL1-0863 & Williams (2018). Creating and integrating a FLOSS product into UK law enforcement. \textit{IFIP Advances in Information and Communication Technology}. DOI: \url{10.1007/978-3-319-92375-8_10}. & 7 \\
RSL1-0906 & Lunhol et al. (2026). INTERACTIVE PLATFORM "VIRTUAL INVESTIGATIVE CASE" IN THE FORMATION OF PROFESSIONAL COMPETENCIES OF FUTURE LAW ENFORCEMENT OFFICERS. \textit{INFORMATION TECHNOLOGIES AND LEARNING TOOLS}. DOI: \url{10.33407/itlt.v114i4.6549}. & 10 \\
RSL1-0908 & Botha et al. (2026). Evaluating an Investigative Process for Cryptocurrency-Related Crimes. \textit{PROCEEDINGS OF THE 21ST INTERNATIONAL CONFERENCE ON CYBER WARFARE AND SECURITY, ICCWS 2026}. & 7 \\
SB1-0104 & Bergman e Popov (2023). Exploring Dark Web Crawlers: A Systematic Literature Review of Dark Web Crawlers and Their Implementation. \textit{IEEE Access}. DOI: \url{10.1109/access.2023.3255165}. & 10 \\
SB1-0202 & Malinverni et al. (2024). SIGNIFICANCE deep learning based platform to fight illicit trafficking of Cultural Heritage goods. \textit{Scientific Reports}. DOI: \url{10.1038/s41598-024-65885-6}. & 7 \\
SB1-0272 & Gaonkar et al. (2026). A Hybrid Approach to Detect Illegal Activities in Dark Web Data. \textit{International Symposium on Digital Forensics and Security}. DOI: \url{10.1109/isdfs69419.2026.11459064}. & 8 \\
SB1-0297 & Sandagiri et al. (2020). ANN Based Crime Detection and Prediction using Twitter Posts and Weather Data. \textit{2020 International Conference on Data Analytics for Business and Industry: Way Towards a Sustainable Economy (ICDABI)}. DOI: \url{10.1109/icdabi51230.2020.9325660}. & 6 \\
SB1-0298 & Yang et al. (2019). Anonymous market product classification based on deep learning. \textit{International Conferences on Artificial Intelligence, Information Processing and Cloud Computing}. DOI: \url{10.1145/3371425.3371467}. & 8 \\
SB1-0306 & A et al. (2026). Automated Dark Web Forensic Intelligence Framework for Cybercrime Investigation Using AI and NLP. \textit{International Conferences on Information Science and System}. DOI: \url{10.1109/iciss67859.2026.11454033}. & 6 \\
SB1-0367 & Tran et al. (2026). Detection of Illicit Content on Online Marketplaces using Large Language Models. \textit{ICON}. DOI: \url{10.1109/icnlp69856.2026.11527538}. & 10 \\
SB1-0404 & Al-harbi e Alghieth (2025). Fine-Tuning Arabic and Multilingual BERT Models for Crime Classification to Support Law Enforcement and Crime Prevention. \textit{International Journal of Advanced Computer Science and Applications}. DOI: \url{10.14569/ijacsa.2025.0160544}. & 10 \\
SB1-0421 & Chatzimarkaki et al. (2023). Harvesting Large Textual and Multimedia Data to Detect Illegal Activities on Dark Web Marketplaces. \textit{BigData Congress [Services Society]}. DOI: \url{10.1109/bigdata59044.2023.10386326}. & 7 \\
SB1-0436 & Ehab et al. (2025). Improving Tor Crawlers Performance Using Deep Learning Methods to bypass CAPTCHA. \textit{2025 IEEE 7th International Conference on Cybernetics, Cognition and Machine Learning Applications (ICCCMLA)}. DOI: \url{10.1109/icccmla66092.2025.11581085}. & 7 \\
SB1-0444 & Kalpakis et al. (2016). Interactive Discovery and Retrieval of Web Resources Containing Home Made Explosive Recipes. \textit{Interacción}. DOI: \url{10.1007/978-3-319-39381-0_20}. & 10 \\
SB1-0465 & Joshi (2024). Multidimensional Forensic Investigation of Onion Sites Based on Fuzzy Encoded LSTM. \textit{Journal of Electrical Systems}. DOI: \url{10.52783/jes.865}. & 6 \\
SB1-0528 & Bergman e Popov (2022). The Digital Detective's Discourse - A toolset for forensically sound collaborative dark web content annotation and collection. \textit{Journal of Digital Forensics, Security and Law}. DOI: \url{10.15394/jdfsl.2022.1740}. & 10 \\
SB1-0547 & Sangher et al. (2023). Towards Safe Cyber Practices: Developing a Proactive Cyber-Threat Intelligence System for Dark Web Forum Content by Identifying Cybercrimes. \textit{Inf.}. DOI: \url{10.3390/info14060349}. & 9 \\
SB1-0584 & Yang et al. (2024). A Reinforcement Learning Approach Combined With Scope Loss Function for Crime Prediction on Twitter (X). \textit{IEEE Access}. DOI: \url{10.1109/access.2024.3473296}. & 9 \\
SB1-0620 & He et al. (2019). Classification of Illegal Activities on the Dark Web. \textit{Proceedings of the 2019 2nd International Conference on Information Science and Systems}. DOI: \url{10.1145/3322645.3322691}. & 10 \\
SB1-0629 & Alkhatib e Basheer (2019). Crawling the Dark Web: A Conceptual Perspective, Challenges and Implementation. \textit{Journal of Digital Information Management}. DOI: \url{10.6025/jdim/2019/17/2/51-60}. & 7 \\
SB1-0633 & Das e Das (2017). Crime analysis against women from online newspaper reports and an approach to apply it in dynamic environment. \textit{2017 International Conference on Big Data Analytics and Computational Intelligence (ICBDAC)}. DOI: \url{10.1109/icbdaci.2017.8070855}. & 10 \\
SB1-0668 & Zhang et al. (2020). dStyle-GAN: Generative Adversarial Network based on Writing and Photography Styles for Drug Identification in Darknet Markets. \textit{Asia-Pacific Computer Systems Architecture Conference}. DOI: \url{10.1145/3427228.3427603}. & 10 \\
SB1-0678 & Cascavilla et al. (2024). Few Images, Many Insights: Illicit Content Detection Using a Limited Number of Images. \textit{ACM Transactions on Intelligent Systems and Technology}. DOI: \url{10.1145/3696458}. & 10 \\
SB1-0689 & Shah et al. (2020). Illicit Activity Detection in Large-Scale Dark and Opaque Web Social Networks. \textit{2020 IEEE International Conference on Big Data (Big Data)}. DOI: \url{10.1109/bigdata50022.2020.9378229}. & 10 \\
SB1-0724 & Porlou et al. (2024). Optimizing an LLM Prompt for Accurate Data Extraction from Firearm-Related Listings in Dark Web Marketplaces. \textit{BigData Congress [Services Society]}. DOI: \url{10.1109/bigdata62323.2024.10825446}. & 9 \\
SB1-0767 & Al-Nabki et al. (2019). ToRank: Identifying the most influential suspicious domains in the Tor network. \textit{Expert systems with applications}. DOI: \url{10.1016/j.eswa.2019.01.029}. & 10 \\
SB1-0781 & Saxena et al. (2023). VendorLink: An NLP approach for Identifying \& Linking Vendor Migrants \& Potential Aliases on Darknet Markets. \textit{Annual Meeting of the Association for Computational Linguistics}. DOI: \url{10.18653/v1/2023.acl-long.481}. & 11 \\
SB1-0827 & Fidalgo et al. (2019). Classifying suspicious content in tor darknet through Semantic Attention Keypoint Filtering. \textit{Digital Investigation. The International Journal of Digital Forensics and Incident Response}. DOI: \url{10.1016/j.diin.2019.05.004}. & 10 \\
SB1-0828 & Al-Nabki et al. (2024). Classifying the content of online notepad services using active learning. \textit{Journal of Intelligence and Information Systems}. DOI: \url{10.1007/s10844-024-00902-8}. & 10 \\
SB1-0837 & Makrynioti et al. (2025). CYBERPATROL: An Intelligent Platform for Monitoring Illicit Online Firearms Trafficking. \textit{BigData Congress [Services Society]}. DOI: \url{10.1109/bigdata66926.2025.11402295}. & 9 \\
SB1-0867 & Nouh e Nurse (2015). Identifying Key-Players in Online Activist Groups on the Facebook Social Network. \textit{2015 IEEE International Conference on Data Mining Workshop (ICDMW)}. DOI: \url{10.1109/icdmw.2015.88}. & 9 \\
SB1-0873 & Al-Nabki et al. (2020). Improving named entity recognition in noisy user-generated text with local distance neighbor feature. \textit{Neurocomputing}. DOI: \url{10.1016/j.neucom.2019.11.072}. & 10 \\
SB1-0877 & Cilleruelo et al. (2020). Interconnection Between Darknets. \textit{IEEE Internet Computing}. DOI: \url{10.1109/mic.2020.3037723}. & 10 \\
SB1-0914 & Al-Nabki et al. (2023). Supervised ranking approach to identify infLuential websites in the darknet. \textit{Applied intelligence (Boston)}. DOI: \url{10.1007/s10489-023-04671-9}. & 10 \\
SB1-1001 & de?Marcos e Domínguez?Díaz (2025). LLM-Based Topic Modeling for Dark Web Q\&A Forums: A Comparative Analysis With Traditional Methods. \textit{IEEE Access}. DOI: \url{10.1109/access.2025.3560543}. & 10 \\
SB1-1025 & Joshi et al. (2019). SummCoder: An unsupervised framework for extractive text summarization based on deep auto-encoders. \textit{Expert systems with applications}. DOI: \url{10.1016/j.eswa.2019.03.045}. & 10 \\
SB1-1026 & Owenson et al. (2018). The darknet's smaller than we thought: The life cycle of Tor Hidden Services. \textit{Digital Investigation. The International Journal of Digital Forensics and Incident Response}. DOI: \url{10.1016/j.diin.2018.09.005}. & 10 \\
\hline
\end{longtable}
\end{small}

\section{RSL 2}
\begin{small}
\begin{longtable}{lp{11.2cm}r}
\hline
\textbf{ID} & \textbf{Referência} & \textbf{Nota} \\ \hline\endhead
RSL2-0065 & Hein et al. (2026). Causal mapping of the risks of using generative AI in software development. \textit{Information and Software Technology}. DOI: \url{10.1016/j.infsof.2026.108202}. & 9 \\
RSL2-0111 & Sergeyuk et al. (2026). Evolving with AI: A Longitudinal Analysis of Developer Logs. \textit{Proceedings - 2026 IEEE/ACM 48th International Conference on Software Engineering, ICSE 2026}. DOI: \url{10.1145/3744916.3787811}. & 10 \\
RSL2-0120 & Guglielmi et al. (2027). How do copilot suggestions impact developers’ frustration and productivity?. \textit{Empirical Software Engineering}. DOI: \url{10.1007/s10664-026-10965-6}. & 9 \\
RSL2-0150 & Ziftci et al. (2026). Catching Developers in the Flow: Low-Latency Agentic Program Repair at Google Scale. \textit{Proceedings - 2026 Proceedings of the 41st IEEE/ACM International Conference on Automated Software Engineering, ASE 2026}. DOI: \url{10.1145/3832783.3834511}. & 8 \\
RSL2-0154 & Ronanki et al. (2026). Recommendations for efficient and responsible LLM adoption within industrial software development. \textit{Information and Software Technology}. DOI: \url{10.1016/j.infsof.2026.108171}. & 9 \\
RSL2-0180 & Zhou et al. (2026). Cognitive Biases in LLM-Assisted Software Development. \textit{Proceedings - 2026 IEEE/ACM 48th International Conference on Software Engineering, ICSE 2026}. DOI: \url{10.1145/3744916.3773104}. & 10 \\
RSL2-0196 & Miller et al. (2026). "maybe We Need Some More Examples:" Individual and Team Drivers of Developer GenAI Tool Use. \textit{Proceedings - 2026 IEEE/ACM 48th International Conference on Software Engineering, ICSE 2026}. DOI: \url{10.1145/3744916.3773155}. & 10 \\
RSL2-0201 & Bolfe et al. (2027). Empirical insights on the use of large language models in software verification and validation: evidence from the Brazilian IT industry. \textit{Empirical Software Engineering}. DOI: \url{10.1007/s10664-026-10971-8}. & 9 \\
RSL2-0205 & Liu e McKinnon (2026). Beyond the Hype: The Efficiency-Throughput Gap with GitHub Copilot. \textit{Communications of the ACM}. DOI: \url{10.1145/3797488}. & 11 \\
RSL2-0214 & Alakmeh et al. (2026). An Eye for AI: Eye-Tracking the Micro-Interruptions of GenAI Code Suggestions. \textit{Proceedings - 2026 IEEE/ACM 48th International Conference on Software Engineering, ICSE 2026}. DOI: \url{10.1145/3744916.3773132}. & 9 \\
RSL2-0216 & Kini et al. (2026). Customizing an LLM for Enterprise Software Engineering. \textit{Proceedings - 2026 Proceedings of the 41st IEEE/ACM International Conference on Automated Software Engineering, ASE 2026}. DOI: \url{10.1145/3832783.3834508}. & 9 \\
RSL2-0225 & Cha et al. (2026). "game Changer" or "overenthusiastic Drunk Acquaintance"? Generative AI Use by Blind and Low Vision Software Professionals in the Workplace. \textit{Proceedings - 2026 IEEE/ACM 48th International Conference on Software Engineering, ICSE 2026}. DOI: \url{10.1145/3744916.3773172}. & 10 \\
RSL2-0240 & Alaswad et al. (2026). Hybrid intelligence effort for software effort estimation in LLM assisted development. \textit{Discover Computing}. DOI: \url{10.1007/s10791-026-10331-6}. & 9 \\
RSL2-0246 & Tabarsi et al. (2027). LLMs’ reshaping of people, processes, products, and society in software development: a qualitative exploration with early adopters. \textit{Empirical Software Engineering}. DOI: \url{10.1007/s10664-026-10917-0}. & 10 \\
RSL2-0322 & Cândido et al. (2027). Developers’ perceptions of automated code review using a generative AI agent: The case of Gandalf. \textit{Journal of Systems and Software}. DOI: \url{10.1016/j.jss.2026.113069}. & 10 \\
RSL2-0378 & Pan et al. (2026). SAINT: Service-level Integration Test Generation with Program Analysis and LLM-based Agents. \textit{Proceedings - 2026 IEEE/ACM 48th International Conference on Software Engineering, ICSE 2026}. DOI: \url{10.1145/3744916.3773185}. & 10 \\
RSL2-0383 & Charoenwet et al. (2026). AgenticSCR: An Autonomous Agentic Secure Code Review for Immature Vulnerabilities Detection. \textit{Proceedings - 2026 Proceedings of the 41st IEEE/ACM International Conference on Automated Software Engineering, ASE 2026}. DOI: \url{10.1145/3832783.3834523}. & 11 \\
RSL2-0393 & Borg et al. (2026). Echoes of AI: Investigating the downstream effects of AI assistants on software maintainability. \textit{Empirical Software Engineering}. DOI: \url{10.1007/s10664-026-10889-1}. & 11 \\
RSL2-0487 & Villani e Kellogg (2026). Command Coding: A Methodological Framework for AI-Accelerated Software Development. \textit{Proceedings - 2026 Proceedings of the 41st IEEE/ACM International Conference on Automated Software Engineering, ASE 2026}. DOI: \url{10.1145/3832783.3834536}. & 10 \\
RSL2-0488 & Pizarro et al. (2026). Beyond the hype: Enabling informed LLM adoption in industry through systematic evaluation. \textit{Science of Computer Programming}. DOI: \url{10.1016/j.scico.2026.103530}. & 11 \\
RSL2-0559 & Martin-Lopez et al. (2026). More Code, Less Understanding? on the Impact of AI Assistants on Developers' Productivity and Code Ownership. \textit{IEEE Transactions on Software Engineering}. DOI: \url{10.1109/TSE.2026.3679627}. & 10 \\
RSL2-0613 & Grandel et al. (2026). ComCat: Expertise-Guided Context Generation to Enhance Code Comprehension. \textit{ACM Transactions on Software Engineering and Methodology}. DOI: \url{10.1145/3742475}. & 10 \\
RSL2-0657 & Kashif et al. (2026). On Developers’ Self-Declaration of AI-Generated Code: An Analysis of Practices. \textit{ACM Transactions on Software Engineering and Methodology}. DOI: \url{10.1145/3771937}. & 10 \\
RSL2-0729 & Zhong et al. (2026). BePilot: An AI Programming Assistant for Compiler Backend Development. \textit{ACM Transactions on Software Engineering and Methodology}. DOI: \url{10.1145/3764585}. & 10 \\
RSL2-0750 & Chen et al. (2026). Programmers Who Use Screen Readers in the Vibe Coding Era: Adaptation, Empowerment, and New Accessibility Landscape. \textit{Conference on Human Factors in Computing Systems - Proceedings }. DOI: \url{10.1145/3772318.3790726}. & 10 \\
RSL2-0761 & Winckler et al. (2026). AI-Assisted Collaboration: Exploring Developer Experience with GitHub Copilot and Windsurf. \textit{IEEE Software}. DOI: \url{10.1109/MS.2026.3671677}. & 9 \\
RSL2-0772 & Fan et al. (2026). When Help Hurts: Verification Load and Fatigue with AI Coding Assistants. \textit{Conference on Human Factors in Computing Systems - Proceedings }. DOI: \url{10.1145/3772318.3791176}. & 10 \\
RSL2-0792 & Yard{\i}m et al. (2026). "the AI tool can't make it any worse." Investigating Developers' Security Behavior with AI Assistants in a Password Storage Study. \textit{Conference on Human Factors in Computing Systems - Proceedings }. DOI: \url{10.1145/3772318.3791693}. & 11 \\
RSL2-0820 & Jang et al. (2026). Evolving Enactions of Expertise: Software Engineers' Evaluation and Demonstration of Coding Expertise with AI Coding Assistants. \textit{Conference on Human Factors in Computing Systems - Proceedings }. DOI: \url{10.1145/3772318.3791260}. & 9 \\
RSL2-0826 & Chang e Lan-gui (2026). The impact of generative artificial intelligence on software development performance: a study based on the theory of technology dominance. \textit{Information and Management}. DOI: \url{10.1016/j.im.2026.104306}. & 8 \\
RSL2-0906 & Shao e Ishengoma (2026). Empirical analysis of generative AI tool adoption in software development. \textit{Information and Software Technology}. DOI: \url{10.1016/j.infsof.2026.108036}. & 6 \\
RSL2-0971 & Tholander e Jönsson (2026). Vibe Coding Entanglements - Repositioning Boundaries of Intention, Authorship, and Responsibility in Programming with Generative AI. \textit{Conference on Human Factors in Computing Systems - Proceedings }. DOI: \url{10.1145/3772318.3791847}. & 10 \\
RSL2-1040 & Alzahrani (2026). Trust and Hallucinations: A Study of 39 Experts on AI-Assisted Requirements Reverse Engineering. \textit{International Journal of Advanced Computer Science and Applications}. DOI: \url{10.14569/IJACSA.2026.0170413}. & 5 \\
RSL2-1129 & Montes e Khojah (2025). Emotional Strain and Frustration in LLM Interactions in Software Engineering. \textit{Proceedings of the 29th International Conference on Evaluation and Assessment in Software Engineering , EASE, 2025 edition, EASE 2025}. DOI: \url{10.1145/3756681.3756951}. & 8 \\
RSL2-1181 & Geruslu et al. (2026). How Software Engineers Engage with AI: A Pragmatic Workflow. \textit{IEEE Software}. DOI: \url{10.1109/MS.2026.3693964}. & 7 \\
RSL2-1182 & Paz-Chen et al. (2026). Google it or ask ChatGPT? Impact of software developers' information search practices on organisational information security. \textit{Behaviour and Information Technology}. DOI: \url{10.1080/0144929X.2026.2654495}. & 11 \\
RSL2-1186 & Sun e Staron (2026). Agentic Pipelines in Embedded Software Engineering: Emerging Practices and Challenges. \textit{Proceedings - 2026 IEEE International Conference on Software Analysis, Evolution and Reengineering, SANER 2026}. DOI: \url{10.1109/SANER67736.2026.00024}. & 7 \\
RSL2-1246 & Aribe e Labastida (2026). The Vibe Shift in Software Engineering: Evaluating AI-Led Conversational Programming for Performance, Cognition, and Responsible Adoption. \textit{International Journal on Advanced Science, Engineering and Information Technology}. DOI: \url{10.18517/IJASEIT.16.4.21812}. & 8 \\
RSL2-1304 & Li et al. (2026). Do comments and expertise still matter? An experiment on programmers’ adoption of AI-generated JavaScript code. \textit{Journal of Systems and Software}. DOI: \url{10.1016/j.jss.2025.112634}. & 11 \\
RSL2-1320 & Nouri et al. (2025). On Simulation-Guided LLM-based Code Generation for Safe Autonomous Driving Software. \textit{Proceedings of the 29th International Conference on Evaluation and Assessment in Software Engineering , EASE, 2025 edition, EASE 2025}. DOI: \url{10.1145/3756681.3756987}. & 9 \\
RSL2-1328 & Oh et al. (2026). How Developers Respond to Insecure AI Code Suggestions under Poisoning Attacks: An Empirical Study of Interface-Level Countermeasures. \textit{IEEE Transactions on Software Engineering}. DOI: \url{10.1109/TSE.2026.3731966}. & 10 \\
RSL2-1385 & Zacharias et al. (2026). Developers’ Dilemma: Opportunities and Pitfalls of Generative AI for Software Development. \textit{Business and Information Systems Engineering}. DOI: \url{10.1007/s12599-026-00998-y}. & 9 \\
RSL2-1500 & Mozannar e Chen (2025). The RealHumanEval: Evaluating Large Language Models’ Abilities to Support Programmers. \textit{Transactions on Machine Learning Research}. & 10 \\
RSL2-1546 & Li et al. (2025). Prompt-with-Me: In-IDE Structured Prompt Management for LLM-Driven Software Engineering. \textit{Proceedings - 2025 40th IEEE/ACM International Conference on Automated Software Engineering, ASE 2025}. DOI: \url{10.1109/ASE63991.2025.00276}. & 9 \\
RSL2-1593 & Zamfirescu-Pereira et al. (2025). Beyond Code Generation: LLM-supported Exploration of the Program Design Space. \textit{Conference on Human Factors in Computing Systems - Proceedings }. DOI: \url{10.1145/3706598.3714154}. & 10 \\
RSL2-1595 & Jensen et al. (2025). Managing expectations towards AI tools for software development: a multiple-case study. \textit{Information Systems and e-Business Management}. DOI: \url{10.1007/s10257-025-00704-7}. & 10 \\
RSL2-1668 & Banh et al. (2025). Copiloting the future: How generative AI transforms Software Engineering. \textit{Information and Software Technology}. DOI: \url{10.1016/j.infsof.2025.107751}. & 8 \\
RSL2-1723 & Ramasamy et al. (2025). AI support for data scientists: An empirical study on workflow and alternative code recommendations. \textit{Empirical Software Engineering}. DOI: \url{10.1007/s10664-025-10622-4}. & 10 \\
RSL2-1753 & Vasconcelos et al. (2025). Generation Probabilities Are Not Enough: Uncertainty Highlighting in AI Code Completions. \textit{ACM Transactions on Computer-Human Interaction}. DOI: \url{10.1145/3702320}. & 10 \\
RSL2-1789 & Sergeyuk et al. (2025). Using AI-based coding assistants in practice: State of affairs, perceptions, and ways forward. \textit{Information and Software Technology}. DOI: \url{10.1016/j.infsof.2024.107610}. & 9 \\
RSL2-1799 & Serafini et al. (2025). Exploring the Impact of Intervention Methods on Developers' Security Behavior in a Manipulated ChatGPT Study. \textit{Conference on Human Factors in Computing Systems - Proceedings }. DOI: \url{10.1145/3706598.3713989}. & 11 \\
RSL2-1840 & Siano et al. (2025). Translating code with Large Language Models and human-in-the-loop feedback. \textit{Information and Software Technology}. DOI: \url{10.1016/j.infsof.2025.107785}. & 9 \\
RSL2-1842 & Cruz et al. (2025). Redefining the Programmer: Human-AI Collaboration, LLMs, and Security in Modern Software Engineering. \textit{Computers, Materials and Continua}. DOI: \url{10.32604/cmc.2025.068137}. & 6 \\
RSL2-1849 & Kirchner et al. (2025). Generative AI Meets Knowledge Management: Insights From Software Development Practices. \textit{Knowledge and Process Management}. DOI: \url{10.1002/kpm.70004}. & 8 \\
RSL2-1878 & Mariotto et al. (2025). From Assessment to Enhancement of Pull Requests at Scale: Aligning Code Reviews with Developer Competencies Using Large Language Models. \textit{International Symposium on Empirical Software Engineering and Measurement}. DOI: \url{10.1109/ESEM64174.2025.00061}. & 10 \\
RSL2-1940 & Flores-Saviaga et al. (2025). The Impact of Generative AI Coding Assistants on Developers Who Are Visually Impaired. \textit{Conference on Human Factors in Computing Systems - Proceedings }. DOI: \url{10.1145/3706598.3714008}. & 8 \\
RSL2-1982 & Keshireddy (2025). Low-Code Development Enhancement Integrating Large Language Models for Intelligent Code Assistance in Oracle APEX. \textit{Indian Journal of Information Sources and Services}. DOI: \url{10.51983/ijiss-2025.IJISS.15.2.46}. & 5 \\
RSL2-2068 & Yang et al. (2025). AskGraph: A Dependency-Aware Code Assistant Powered by Code Graphs and LLM-Generated Cypher Queries. \textit{Proceedings - 2025 IEEE International Conference on Software Maintenance and Evolution, ICSME 2025}. DOI: \url{10.1109/ICSME64153.2025.00065}. & 10 \\
RSL2-2216 & Erhabor et al. (2025). Measuring the Runtime Performance of C++ Code Written by Humans Using Github Copilot. \textit{Proceedings - International Conference on Software Engineering}. DOI: \url{10.1109/ICSE55347.2025.00059}. & 9 \\
RSL2-2436 & Di e Zhang (2025). Enhancing Code Generation via Bidirectional Comment-Level Mutual Grounding. \textit{Proceedings - International Conference on Software Engineering}. DOI: \url{10.1109/ICSE55347.2025.00165}. & 11 \\
RSL2-2477 & Klemmer et al. (2024). Using AI Assistants in Software Development: A Qualitative Study on Security Practices and Concerns*. \textit{CCS 2024 - Proceedings of the 2024 ACM SIGSAC Conference on Computer and Communications Security}. DOI: \url{10.1145/3658644.3690283}. & 10 \\
RSL2-2484 & Tufano et al. (2025). Deep Learning-based Code Reviews: A Paradigm Shift or a Double-Edged Sword?. \textit{Proceedings - International Conference on Software Engineering}. DOI: \url{10.1109/ICSE55347.2025.00060}. & 11 \\
RSL2-2485 & Kumar et al. (2025). Why AI Agents Still Need You: Findings from Developer-Agent Collaborations in the Wild. \textit{Proceedings - 2025 40th IEEE/ACM International Conference on Automated Software Engineering, ASE 2025}. DOI: \url{10.1109/ASE63991.2025.00043}. & 10 \\
RSL2-2566 & Liang et al. (2024). A Large-Scale Survey on the Usability of AI Programming Assistants: Successes and Challenges. \textit{Proceedings - International Conference on Software Engineering}. DOI: \url{10.1145/3597503.3608128}. & 10 \\
RSL2-2576 & Perry et al. (2023). Do Users Write More Insecure Code with AI Assistants?. \textit{CCS 2023 - Proceedings of the 2023 ACM SIGSAC Conference on Computer and Communications Security}. DOI: \url{10.1145/3576915.3623157}. & 11 \\
RSL2-2593 & Mozannar et al. (2024). Reading Between the Lines: Modeling User Behavior and Costs in AI-Assisted Programming. \textit{Conference on Human Factors in Computing Systems - Proceedings}. DOI: \url{10.1145/3613904.3641936}. & 10 \\
RSL2-2624 & Russo (2024). Navigating the Complexity of Generative AI Adoption in Software Engineering. \textit{ACM Transactions on Software Engineering and Methodology}. DOI: \url{10.1145/3652154}. & 9 \\
RSL2-2650 & Weber et al. (2024). Significant Productivity Gains through Programming with Large Language Models. \textit{Proceedings of the ACM on Human-Computer Interaction}. DOI: \url{10.1145/3661145}. & 10 \\
RSL2-2710 & Gu et al. (2024). How Do Analysts Understand and Verify AI-Assisted Data Analyses?. \textit{Conference on Human Factors in Computing Systems - Proceedings}. DOI: \url{10.1145/3613904.3642497}. & 9 \\
RSL2-2746 & Barke et al. (2023). Grounded Copilot: How Programmers Interact with Code-Generating Models. \textit{Proceedings of the ACM on Programming Languages}. DOI: \url{10.1145/3586030}. & 8 \\
RSL2-2812 & Asare et al. (2024). A User-Centered Security Evaluation of Copilot. \textit{Proceedings - International Conference on Software Engineering}. DOI: \url{10.1145/3597503.3639154}. & 10 \\
RSL2-2815 & Ferdowsi et al. (2024). Validating AI-Generated Code with Live Programming. \textit{Conference on Human Factors in Computing Systems - Proceedings}. DOI: \url{10.1145/3613904.3642495}. & 10 \\
RSL2-2839 & Kuhail et al. (2024). “Will I be replaced?” Assessing ChatGPT's effect on software development and programmer perceptions of AI tools. \textit{Science of Computer Programming}. DOI: \url{10.1016/j.scico.2024.103111}. & 10 \\
RSL2-2843 & Oh et al. (2024). Poisoned ChatGPT Finds Work for Idle Hands: Exploring Developers' Coding Practices with Insecure Suggestions from Poisoned AI Models. \textit{Proceedings - IEEE Symposium on Security and Privacy}. DOI: \url{10.1109/SP54263.2024.00046}. & 11 \\
RSL2-2903 & Spain e Pashchenko (2023). Early Formalization of AI-tools Usage in Software Engineering in Europe: Study of 2023. \textit{International Journal of Information Technology and Computer Science}. DOI: \url{10.5815/ijitcs.2023.06.03}. & 5 \\
RSL2-3300 & Ziegler et al. (2024). Measuring GitHub Copilot's Impact on Productivity. \textit{COMMUNICATIONS OF THE ACM}. DOI: \url{10.1145/3633453}. & 9 \\
RSL2-3381 & Choudhuri et al. (2026). To Copilot and Beyond: 22 AI Systems Developers Want Built. \textit{2026 IEEE International Conference on Software Maintenance and Evolution (ICSME)}. DOI: \url{10.1109/ICSME67971.2026.00060}. & 8 \\
RSL2-3387 & Richards et al. (2026). No Two Developers Think Alike: How Problem-Solving Styles and Experience Shape Needs in Conversational Interaction with Copilot. \textit{2026 IEEE International Conference on Software Maintenance and Evolution (ICSME)}. DOI: \url{10.1109/ICSME67971.2026.00022}. & 10 \\
RSL2-3693 & Olewicki et al. (2026). Impact of an LLM-based Review Assistant in Practice: A Mixed Open-/Closed-source Case Study. \textit{IEEE Transactions on Software Engineering}. DOI: \url{10.1109/TSE.2026.3663093}. & 10 \\
RSL2-3733 & Bull e Kharrufa (2024). Generative Artificial Intelligence Assistants in Software Development Education: A Vision for Integrating Generative Artificial Intelligence Into Educational Practice, Not Instinctively Defending Against It. \textit{IEEE Software}. DOI: \url{10.1109/MS.2023.3300574}. & 9 \\
RSL2-3768 & Ullrich et al. (2025). From Requirements to Code: Understanding Developer Practices in LLM-Assisted Software Engineering. \textit{2025 IEEE 33rd International Requirements Engineering Conference (RE)}. DOI: \url{10.1109/RE63999.2025.00032}. & 9 \\
RSL2-3833 & Shin et al. (2025). Prompt Engineering or Fine-Tuning: An Empirical Assessment of LLMs for Code. \textit{2025 IEEE/ACM 22nd International Conference on Mining Software Repositories (MSR)}. DOI: \url{10.1109/MSR66628.2025.00082}. & 11 \\
RSL2-3904 & Fakhoury et al. (2024). LLM-Based Test-Driven Interactive Code Generation: User Study and Empirical Evaluation. \textit{IEEE Transactions on Software Engineering}. DOI: \url{10.1109/TSE.2024.3428972}. & 11 \\
RSL2-3990 & Tie et al. (2026). "Should I Give Up Now?" Investigating LLM Pitfalls in Software Engineering. \textit{ACM Trans. Softw. Eng. Methodol.}. DOI: \url{10.1145/3801972}. & 9 \\
RSL2-4029 & Moreno-Lumbreras et al. (2026). BonsAIDE: An Extended Vision for Human–AI Interaction in IDEs. \textit{ACM Trans. Softw. Eng. Methodol.}. DOI: \url{10.1145/3793681}. & 5 \\
RSL2-4034 & Brandebusemeyer et al. (2026). Developers’ Experience with Generative AI Beyond Productivity Assessment – Insights from an Empirical Mixed-Methods Field Study. \textit{ACM Trans. Softw. Eng. Methodol.}. DOI: \url{10.1145/3847122}. & 11 \\
RSL2-4042 & Gao et al. (2026). Evaluating the Impact of Explainable AI on Trust in AI-Assisted Code Review. \textit{Proc. ACM Softw. Eng.}. DOI: \url{10.1145/3832184}. & 10 \\
RSL2-4074 & Vella e Blincoe (2026). The Impact of AI Coding Assistants on Software Engineering: A Longitudinal Study. \textit{ACM Trans. Softw. Eng. Methodol.}. DOI: \url{10.1145/3861343}. & 10 \\
RSL2-4104 & Choudhuri et al. (2026). What Needs Attention? Prioritizing Drivers of Developers’ Trust and Adoption of Generative AI. \textit{ACM Trans. Softw. Eng. Methodol.}. DOI: \url{10.1145/3842739}. & 8 \\
RSL2-4151 & Alami et al. (2026). Engagement in Code Review: Emotional, Behavioral, and Cognitive Dimensions in Peer vs. LLM Interactions. \textit{ACM Trans. Softw. Eng. Methodol.}. DOI: \url{10.1145/3830405}. & 11 \\
SB2-0075 & Kang et al. (2024). A Quantitative and Qualitative Evaluation of LLM-Based Explainable Fault Localization. \textit{Proceedings of the ACM on software engineering.}. DOI: \url{10.1145/3660771}. & 11 \\
SB2-0132 & Alami et al. (2025). Accountability in Code Review: The Role of Intrinsic Drivers and the Impact of LLMs. \textit{ACM Transactions on Software Engineering and Methodology}. DOI: \url{10.1145/3721127}. & 11 \\
SB2-0162 & Maddila et al. (2026). Agentic Program Repair From Test Failures at Scale: A Neuro-Symbolic Approach With Static Analysis and Test Execution Feedback. \textit{IEEE Transactions on Software Engineering}. DOI: \url{10.1109/tse.2026.3696849}. & 10 \\
SB2-0206 & Murali et al. (2024). AI-Assisted Code Authoring at Scale: Fine-Tuning, Deploying, and Mixed Methods Evaluation. \textit{Proceedings of the ACM on software engineering.}. DOI: \url{10.1145/3643774}. & 11 \\
SB2-0207 & Isong e Vhahangwele (2026). AI-Assisted Development Tools and Team Dynamics in South African Software Engineering Teams. \textit{Journal of Information Systems and Informatics}. DOI: \url{10.63158/journalisi.v8i3.1615}. & 8 \\
SB2-0258 & Inoue et al. (2026). An Empirical Investigation on Source Code Comprehension Performance Between Naming and Structural Issues With and Without Generative AI Assistance. \textit{IEEE Transactions on Software Engineering}. DOI: \url{10.1109/tse.2026.3702933}. & 10 \\
SB2-0279 & Quiles et al. (2026). An intelligent low-code platform for building task-oriented LLM-based chatbots. \textit{Journal of Computer Languages}. DOI: \url{10.1016/j.cola.2026.101401}. & 9 \\
SB2-0308 & Costa et al. (2025). Applying generative artificial intelligence for vulnerability fixing in a proprietary software ecosystem. \textit{Journal of Systems and Software}. DOI: \url{10.1016/j.jss.2025.112723}. & 11 \\
SB2-0426 & Khojah et al. (2024). Beyond Code Generation: An Observational Study of ChatGPT Usage in Software Engineering Practice. \textit{Proceedings of the ACM on software engineering.}. DOI: \url{10.1145/3660788}. & 10 \\
SB2-0471 & Kruse et al. (2024). Can Developers Prompt? A Controlled Experiment for Code Documentation Generation. \textit{IEEE International Conference on Software Maintenance (ICSM)}. DOI: \url{10.1109/icsme58944.2024.00058}. & 11 \\
SB2-0554 & Kim et al. (2026). CodeVoyager: Integrating Interactive Visual Aids with LLMs for Code Comprehension. \textit{ACM International Conference on Intelligent User Interfaces (IUI)}. DOI: \url{10.1145/3742413.3789057}. & 11 \\
SB2-0597 & Wadhwa et al. (2024). CORE: Resolving Code Quality Issues using LLMs. \textit{Proceedings of the ACM on software engineering.}. DOI: \url{10.1145/3643762}. & 11 \\
SB2-0652 & Salomon et al. (2026). Deriving Personas to Describe Software Development Team Members in the Era of GenAI. \textit{IEEE International Conference on Software Maintenance (ICSM)}. DOI: \url{10.1109/icsme67971.2026.00042}. & 9 \\
SB2-0678 & Stalnaker et al. (2025). Developer Perspectives on Licensing and Copyright Issues Arising from Generative AI for Software Development. \textit{ACM Transactions on Software Engineering and Methodology}. DOI: \url{10.1145/3743133}. & 10 \\
SB2-0690 & Jiang et al. (2022). Discovering the Syntax and Strategies of Natural Language Programming with Generative Language Models. \textit{CHI Conference on Human Factors in Computing Systems}. DOI: \url{10.1145/3491102.3501870}. & 8 \\
SB2-0844 & Kang et al. (2024). Explainable automated debugging via large language model-driven scientific debugging. \textit{Empirical Software Engineering}. DOI: \url{10.1007/s10664-024-10594-x}. & 10 \\
SB2-0846 & Valenzuela-Toledo et al. (2025). Explaining GitHub Actions Failures with Large Language Models: Challenges, Insights, and Limitations. \textit{IEEE International Conference on Program Comprehension (ICPC)}. DOI: \url{10.1109/icpc66645.2025.00037}. & 9 \\
SB2-0873 & Liu et al. (2026). Exploring the integration of large language models in industrial test maintenance processes. \textit{Journal of Systems and Software}. DOI: \url{10.1016/j.jss.2026.113049}. & 10 \\
SB2-0928 & Salomon et al. (2026). From Disruptions to Discussions: How GenAI Impacts Human Interactions in Software Development. \textit{IEEE Transactions on Software Engineering}. DOI: \url{10.1109/tse.2026.3655626}. & 9 \\
SB2-0933 & Feng et al. (2026). From Junior to Senior: Allocating Agency and Navigating Professional Growth in Agentic AI-Mediated Software Engineering. \textit{CHI Conference on Human Factors in Computing Systems}. DOI: \url{10.1145/3772318.3791642}. & 10 \\
SB2-1143 & Arisantoso et al. (2026). Human-AI Work Redistribution Across the Software Development Life Cycle: A Socio-Technical Study of Indonesian Developers. \textit{Applied Information System and Management (AISM)}. DOI: \url{10.15408/aism.v9i2.53455}. & 9 \\
SB2-1193 & You et al. (2026). Industrial Practice of LLM-Based Test Case Carving and Assertion Generation (Experience Paper). \textit{Proceedings of the ACM on software engineering.}. DOI: \url{10.1145/3832099}. & 8 \\
SB2-1255 & Kabir et al. (2024). Is Stack Overflow Obsolete? An Empirical Study of the Characteristics of ChatGPT Answers to Stack Overflow Questions. \textit{CHI Conference on Human Factors in Computing Systems}. DOI: \url{10.1145/3613904.3642596}. & 11 \\
SB2-1467 & Pomian et al. (2024). Next-Generation Refactoring: Combining LLM Insights and IDE Capabilities for Extract Method. \textit{IEEE International Conference on Software Maintenance (ICSM)}. DOI: \url{10.1109/icsme58944.2024.00034}. & 10 \\
SB2-1525 & Martinovic e Rozic (2025). Perceived Impact of AI-Based Tooling on Software Development Code Quality. \textit{SN Computer Science}. DOI: \url{10.1007/s42979-024-03608-4}. & 7 \\
SB2-1605 & Sun et al. (2026). Quality assurance of LLM-generated code: Addressing non-functional quality characteristics. \textit{Journal of Systems and Software}. DOI: \url{10.1016/j.jss.2026.112885}. & 11 \\
SB2-1636 & Pandini et al. (2026). Refactoring architectural smells with large language models: an exploratory study on cyclic dependencies. \textit{Empirical Software Engineering}. DOI: \url{10.1007/s10664-026-10939-8}. & 11 \\
SB2-1783 & Kemell et al. (2025). Still just personal assistants? – A multiple case study of generative AI adoption in software organizations. \textit{Information and Software Technology}. DOI: \url{10.1016/j.infsof.2025.107805}. & 10 \\
SB2-1811 & Abedu et al. (2026). Synergizing LLMs and Knowledge Graphs: A Novel Approach to Software Repository-Related Question Answering. \textit{ACM Transactions on Software Engineering and Methodology}. DOI: \url{10.1145/3796510}. & 11 \\
SB2-1818 & Qian e Wexler (2024). Take It, Leave It, or Fix It: Measuring Productivity and Trust in Human-AI Collaboration. \textit{ACM International Conference on Intelligent User Interfaces (IUI)}. DOI: \url{10.1145/3640543.3645198}. & 11 \\
SB2-1833 & Pashchenko (2026). Teoría y práctica del desarrollo de software aumentado con IA en Europa en 2025. \textit{Revista Multidisciplinar Epistemología de las Ciencias}. DOI: \url{10.71112/pahqyx95}. Relato complementar de SB2-1960. & 4 \\
SB2-1918 & Fortes et al. (2026). The Productivity Paradox of AI-Powered Development. \textit{IEEE Software}. DOI: \url{10.1109/ms.2026.3658651}. & 10 \\
SB2-1960 & Pashchenko (2025). Theory and Practice of AI-Augmented Software Development in Europe. \textit{Real-World AI Systems}. DOI: \url{10.64797/rwas.v1i2.104}. & 4 \\
SB2-2049 & Mousavi et al. (2026). Understanding the Impact of AI Code Assistants on Security API Usage: An Empirical Study. \textit{Lecture notes in computer science}. DOI: \url{10.1007/978-3-032-42939-1_12}. & 10 \\
SB2-2107 & Mock e Russo (2026). Vibe Coding: An experiment with test-driven development. \textit{Empirical Software Engineering}. DOI: \url{10.1007/s10664-026-10950-z}. & 11 \\
SB2-2163 & Kares et al. (2026). Working with Large Language Models: impacts on performance, efficiency and perceived responsibility. \textit{Behaviour and Information Technology}. DOI: \url{10.1080/0144929x.2026.2666279}. & 10 \\
SB2-2369 & Shen et al. (2024). Debugging Defective Visualizations: Empirical Insights Informing a Human-AI Co?Debugging System. \textit{International Conference on Human Factors in Computing Systems}. DOI: \url{10.1145/3772318.3791441}. & 11 \\
SB2-2374 & Hossain e LaToza (2026). DesignFix: Context-Augmented Repair of Design Rule Violations. \textit{Proceedings of the 41st IEEE/ACM International Conference on Automated Software Engineering}. DOI: \url{10.1145/3832783.3837422}. & 11 \\
SB2-2551 & Liu et al. (2025). MACEDON : Supporting Programmers with Real-Time Multi-Dimensional Code Evaluation and Optimization. \textit{ACM Symposium on User Interface Software and Technology}. DOI: \url{10.1145/3746059.3747655}. & 9 \\
SB2-2613 & Zhou et al. (2026). RESTOR: Automated Test Oracle Generation for RESTful APIs via Reinforcement Learning. \textit{Proceedings of the ACM on Software Engineering}. DOI: \url{10.1145/3832185}. & 11 \\
\hline
\end{longtable}
\end{small}
